\documentclass{article}
\usepackage{amsmath}
\usepackage{graphicx} 
\usepackage{float}
\usepackage{listings}

\begin{document}
	
	\title{Teoria de las Superficies}
	\author{Por: David Duarte}
	\date{\today}
	
	\maketitle
	

	
	En geometría diferencial, las representaciones paramétricas regulares son fundamentales para describir superficies locales de manera precisa. Una superficie simple puede ser definida por una función vectorial $	\textbf{r}(\theta, \phi)$, donde $\theta$ y $\phi$ son parámetros que varían en un cierto dominio. 
	
	\section{Representaciones Paramétricas}
		
		Las representaciones paramétricas regulares son una forma de describir superficies en el espacio tridimensional mediante funciones vectoriales que dependen de parámetros. Estas representaciones son útiles para visualizar y estudiar diferentes tipos de superficies de manera sistemática.
		
		Existen varias formas de parametrizar una superficie, y la elección de los parámetros puede variar según la forma de la superficie y la conveniencia matemática. Algunos ejemplos comunes de representaciones paramétricas regulares incluyen:
		
		1. \textbf{Esfera}: Una esfera de radio $R$ centrada en el origen puede ser parametrizada por las coordenadas esféricas $(\theta, \phi)$ de la siguiente manera:
		
		$$
		\textbf{r}(\theta, \phi) = \begin{pmatrix} R \cos(\theta) \sin(\phi) \\ R \sin(\theta) \sin(\phi) \\ R \cos(\phi) \end{pmatrix}
		$$
		
		donde $0 \leq \theta < 2\pi$ y $0 \leq \phi \leq \pi$.
		
		2. \textbf{Toroide}: Un toroide con radios $R$ y $r$ puede ser parametrizado por los ángulos $\theta$ y $\phi$ de la siguiente manera:
		
		$$
		\textbf{r}(\theta, \phi) = \begin{pmatrix} (R + r \cos(\theta)) \cos(\phi) \\ (R + r \cos(\theta)) \sin(\phi) \\ r \sin(\theta) \end{pmatrix}
		$$
		
		donde $0 \leq \theta < 2\pi$ y $0 \leq \phi < 2\pi$.
		
		3. \textbf{Cilindro}: Un cilindro de radio $R$ centrado en el eje $z$ puede ser parametrizado por la altura $z$ y el ángulo $\theta$ de la siguiente manera:
		
		$$
		\textbf{r}(z, \theta) = \begin{pmatrix} R \cos(\theta) \\ R \sin(\theta) \\ z \end{pmatrix}
		$$
		
		donde $0 \leq z \leq H$ y $0 \leq \theta < 2\pi$.
		
		Estos son solo algunos ejemplos de representaciones paramétricas regulares de superficies comunes. La elección de los parámetros y la forma de la función vectorial $\textbf{r}$ pueden variar dependiendo de la superficie que se esté estudiando.
		
	\subsection{Ejemplos}
	
		A modo de ejemplo podemos visualizar lo siguiente:
		
	
			
			\textbf{Ejemplo 1: Superficie de revolución}
			
			Consideremos la superficie de revolución generada por la curva $y = f(x)$ alrededor del eje $x$. La parametrización de esta superficie puede ser dada por:
			
			$$
			\textbf{r}(x, \theta) = \begin{pmatrix} x \\ f(x) \cos(\theta) \\ f(x) \sin(\theta) \end{pmatrix}
			$$
			
			donde $a \leq x \leq b$ y $0 \leq \theta < 2\pi$. Esta representación paramétrica describe cada punto en la superficie de revolución en función de dos parámetros: la variable $x$ que recorre la curva generatriz y el ángulo $\theta$ que varía alrededor del eje de rotación.
			
			Ver ejemplo de código en Python
			
			\textbf{Ejemplo 2: Superficie hiperbólica}
			
			La superficie hiperbólica puede ser parametrizada de la siguiente manera:
			
			$$
			\textbf{r}(u, v) = \begin{pmatrix} a \cosh(u) \cos(v) \\ a \sinh(u) \sin(v) \\ b \cosh(u) \end{pmatrix}
			$$
			
			donde $-\infty < u < \infty$ y $0 \leq v < 2\pi$. En esta parametrización, los parámetros $u$ y $v$ permiten describir puntos en la superficie hiperbólica en términos de las funciones hiperbólicas coseno hiperbólico ($\cosh$) y seno hiperbólico ($\sinh$).
			
			Estos ejemplos ilustran cómo las representaciones paramétricas pueden utilizarse para describir una variedad de superficies de manera precisa y sistemática.
			
			Ver ejemplo de código en Python
			
		
	
		
		
\section{Plano Tangente}
	
	
	El plano tangente a la superficie en un punto $(u_0, v_0)$ se puede obtener a partir de las derivadas parciales de la función $\textbf{r}(u, v)$. La ecuación del plano tangente es:
	
	$$
	\textbf{r}_u(u_0, v_0) \cdot (x - x_0) + \textbf{r}_v(u_0, v_0) \cdot (y - y_0) + \textbf{r}(u_0, v_0) \cdot (z - z_0) = 0
	$$
	
	donde $\textbf{r}_u$ y $\textbf{r}_v$ son las derivadas parciales de $\textbf{r}$ respecto a $u$ y $v$ respectivamente.
	
	La recta normal a la superficie en un punto $(u_0, v_0)$ es perpendicular al plano tangente y su vector dirección es proporcional al producto cruz de los vectores tangentes $\textbf{r}_u$ y $\textbf{r}_v$:
	
	$$
	\textbf{N}(u_0, v_0) = \frac{\textbf{r}_u \times \textbf{r}_v}{|\textbf{r}_u \times \textbf{r}_v|}
	$$
	
	Las propiedades topológicas de las superficies simples se refieren a características como la conectividad, la orientabilidad y la compacidad de la superficie.
	
	
En el siguiente ejemplo se observa un plano tangente a una superficie en un punto dado:


\begin{figure}[H]
	\centering
	\includegraphics[width=0.7\linewidth]{"C:/Users/magno/Documents/Proyectos de Programación/Python/Superficies_parametricas/Figure_2"}
	\caption{}
	\label{fig:figure1}
\end{figure}

Las ecuaciones paramétricas de la superficie son:

	$$
	\textbf{r}(u, v) = \begin{pmatrix} u \\ v \\ u^3-v^3 \end{pmatrix}
	$$
	
	Y el punto de tangencia es:
	
	$$ \textbf{P}_{t}(u, v) = (1,1) $$
	
\section{Recta normal}

La recta normal a la superficie paramétrica dada por 

$$\boldsymbol{r}(u, v) = (x(u, v), y(u, v), z(u, v))$$

 en el punto $(u_0, v_0)$ está dada por el vector normal $\boldsymbol{n}$ en ese punto. 

El vector normal $\boldsymbol{n}$ se puede calcular como el producto cruz entre los vectores tangentes $\boldsymbol{r}_u$ y $\boldsymbol{r}_v$ en el punto $(u_0, v_0)$, y se normaliza para obtener la dirección de la recta normal.

Tomando como ejemplo, la superficie anterior, podemos calcular la recta normal a dicha superficie, en el mismo punto de tangencia usado para el plano tangente:

\begin{figure}[H]
	\centering
	\includegraphics[width=0.7\linewidth]{"C:/Users/magno/Documents/Proyectos de Programación/Python/Superficies_parametricas/Figure_3"}
	\caption{}
	\label{fig:figure3}
\end{figure}

	
	\section{Anexos}
	
	
\begin{lstlisting}[language=Python, caption={Ejemplo de código Python para superficie de Revolución}, label={python}]
	import numpy as np
	import matplotlib.pyplot as plt
	from mpl_toolkits.mplot3d import Axes3D
	
	
	def surface_of_revolution(f, a, b):
	x = np.linspace(a, b, 100)
	theta = np.linspace(0, 2*np.pi, 100)
	X, Theta = np.meshgrid(x, theta)
	Y = f(X) * np.cos(Theta)
	Z = f(X) * np.sin(Theta)
	return X, Y, Z
	
	def f(x):
	return np.sin(x)
	
	X, Y, Z = surface_of_revolution(f, 0, 2*np.pi)
	
	fig = plt.figure()
	ax = fig.add_subplot(111, projection='3d')
	ax.plot_surface(X, Y, Z, cmap='viridis')
	ax.set_title('Superficie de Revolucion')
	plt.show()
	

\end{lstlisting}

\begin{lstlisting}[language=Python, caption={Ejemplo de código Python para superficie hiperbólica}, label={python}]
	import numpy as np
	import matplotlib.pyplot as plt
	from mpl_toolkits.mplot3d import Axes3D
	
	

	def hyperbolic_surface(a, b):
	u = np.linspace(-5, 5, 100)
	v = np.linspace(0, 2*np.pi, 100)
	U, V = np.meshgrid(u, v)
	X = a * np.cosh(U) * np.cos(V)
	Y = a * np.sinh(U) * np.sin(V)
	Z = b * np.cosh(U)
	return X, Y, Z
	
	X, Y, Z = hyperbolic_surface(1, 2)
	
	fig = plt.figure()
	ax = fig.add_subplot(111, projection='3d')
	ax.plot_surface(X, Y, Z, cmap='plasma')
	ax.set_title('Superficie Hiperbolica')
	plt.show()
\end{lstlisting}
	
\section{Ejercicios Propuestos}

\begin{enumerate}
	\item Escriba un codigo Python que grafique lo siguiente:
	
		\begin{itemize}
			\item Una superficie esférica de radio = 3
			\item Un cilindro de radio r = 1 y altura h = 10
			\item Un toroide de radio R = 3 y radio r = 1
		\end{itemize}
		
	\item Escriba un código Python para determinar y graficar el plano tangente y la recta normal a la superficie :
	
	$$
	\textbf{r}(u, v) = \begin{pmatrix} u \\ v \\ u^3-v^2 \end{pmatrix}
	$$
	
	en el punto $ \textbf{P}(u, v) = (0, 0) $
	
\end{enumerate}


	
\end{document}
